\documentclass[10pt,a4paper]{article}
\usepackage[T1]{fontenc}
\usepackage{lmodern,amsmath,amssymb,microtype}
\usepackage[margin=20mm]{geometry}
\usepackage[hidelinks]{hyperref}
\pdfinfoomitdate=1
\pdftrailerid{}
\pdfsuppressptexinfo=15
\setlength{\parindent}{0pt}
\setlength{\parskip}{5pt}
\begin{document}
\begin{center}
{\Large Eclipse-constrained time GPs and conditional NUTS}\par\vspace{4pt}
\texttt{exojax-escape} --- 25 November 2026
\end{center}
\textbf{Purpose.} Infer a fixed atmospheric spectrum in the presence of
correlated time variation. Temperature, populations, opacity and wind remain
at the saved atmosphere. A shared velocity offset and positive contrast
describe the planet. NUTS samples this conditional model, not the full
thermochemical atmosphere.

\textbf{Eclipse control.} Select science exposures at $0.30\leq\phi<0.435$
and 15 conservative eclipse exposures at $0.47\leq\phi\leq0.5$ in the second night.
Intermediate exposures are excluded because partial stellar occultation is
not implemented. Planet flux is zero in the fully eclipsed exposures,
conditional on the supplied phase convention, circular geometry and emission
confined within $3R_p$. Independent timing/ephemeris metadata are absent.
The 46-minute eclipse segment starts about 79 minutes after the last selected
science exposure. It constrains noise statistics more readily than the
specific earlier realization. The first night has no eclipse control.

\textbf{Mean and covariance.} For each night/order, use native stellar pixels
at $110$--$390\,{\rm km\,s^{-1}}$. The archived log ratio is
$y_{ep}=\log(1+r_{ep})$ and $m_{ep}(\theta)$ is the logarithm of the
continuum-normalized native star-plus-planet prediction:
\begin{equation}
y_{ep}=m_{ep}(\theta)+c_p+a(t_e)+x_p b(t_e)+\epsilon_{ep},\qquad
x_p=(v_{*,p}-250)/140.
\label{eq:mean}
\end{equation}
$c_p$ absorbs an arbitrary common nightly log spectrum. Independent
zero-mean processes $a,b$ describe a time-varying grey term and spectral
slope, with Mat\'ern-$3/2$ covariance
\begin{equation}
k_j(t,t')=A_j^2(1+\sqrt3\,|t-t'|/\ell)
                 \exp(-\sqrt3\,|t-t'|/\ell),\qquad j=0,1.
\label{eq:kernel}
\end{equation}
One length is shared by the two modes within a night/order; amplitudes are
separate. Independent night/order realizations do not assume identical
instrumental histories. This initial model does not recover cross-order
covariance or arbitrary narrow spectral systematics.

The eclipse residuals reveal substantial correlation between original adjacent
pixels. The primary NUTS comparison therefore keeps original even-numbered
pixels; odd-numbered pixels form a separate sensitivity analysis. This fixed
selection uses the eclipse control before the main fit. It loses information
and reduces the clearest neighbor covariance; it does not prove that all
remaining correlations vanish. The two interleaves are never multiplied as
independent observations.

White noise has covariance $s^2D_0$, with the fixed first-order log-error proxy
$D_{0,ep}=[\sigma_{ep}/(1+r_{ep})]^2$. The scale $s$ is inferred: supplied
errors are empirical residual scatter, not an independently known photon
noise floor. Adding GP variance to a fixed empirical total variance could
count variability twice.

\textbf{Remove the common reference from mean and covariance.}
An orthonormal contrast operator satisfies $CC^T=I$ and
$CD_0^{-1/2}X_c=0$, where $X_c$ assigns one constant per native wavelength.
Its rank is $R=\sum_p\max(E_p-1,0)$ for $E_p$ valid exposures. It does not count
all pairwise differences independently. Set $d=CD_0^{-1/2}(y-m)$ and
$G=CD_0^{-1/2}B$, where $B$ maps coefficients $(a_e,b_e)$ onto
 equation~\eqref{eq:mean}. For their covariance $K=LL^T$,
\begin{equation}
\log p(d)=-\tfrac12\left[d^T\Sigma^{-1}d+\log|\Sigma|+R\log(2\pi)\right],
\qquad\Sigma=s^2I+GLL^TG^T.
\label{eq:likelihood}
\end{equation}
GP functions are integrated out analytically. The determinant penalizes
excessive covariance; fitting only a corrected residual norm would not give
this posterior.

\newpage
\textbf{Efficient evaluation.}
With $H=G^TG$, $z=G^Td$, $S=I+L^THL/s^2$, and $u=L^Tz/s^2$,
\begin{equation}
d^T\Sigma^{-1}d=\frac{d^Td}{s^2}-u^TS^{-1}u,\qquad
\log|\Sigma|=2R\log s+\log|S|.
\label{eq:woodbury}
\end{equation}
The largest factorization uses twice the number of exposures in one
night/order. Fixed weighted native sums construct $H,z,d^Td$ directly,
without a pixel-size covariance or explicit large $C$.
A relative $10^{-10}$ diagonal regularization stabilizes the correlation
Cholesky and scales with $A_j^2$. Zero GP amplitude therefore recovers the
white-noise model exactly. Only constants independent of every fitted
parameter are omitted. Differences across sample masks need their measure
constants before being called absolute predictive densities.

\textbf{Declared priors.}
The shared planet has $\Delta v\sim\mathcal U(-15,15)$ in
${\rm km\,s^{-1}}$ and $\log A_{\rm planet}\sim\mathcal N(0,0.7^2)$.
For each night/order,
\begin{equation}
\log A_j\sim\mathcal N(\log(300\times10^{-6}),1),\quad
\log(\ell/{\rm min})\sim\mathcal N(\log20,1),\quad
\log s\sim\mathcal N(0,0.35^2).
\end{equation}
The white-only comparison uses identical planet and white-scale priors and
data. Positive GP amplitude priors exclude zero by construction: an interval
above zero is not a detection. Report prior sensitivity and prediction tests
separately from sampler diagnostics. Times are relative within a night,
obtained from phase and the adopted 1.48-day period; no absolute timing
calibration is recovered.
The GP is evaluated at exposure centers; its kernel does not average over
the 150-second exposure duration.

\textbf{Validation and interpretation.}
Independent dense projected-covariance tests check the likelihood,
determinants, derivatives and conditional moments. A synthetic common master
must cancel. Artificial planet signals with and without correlated noise
test recovery and the risk of absorbing a moving planetary line. Contiguous
eclipse blocks are withheld before fitting their noise hyperparameters.
NUTS convergence, numerical correspondence and noise-model adequacy remain
separate claims.

A GP mean fitted to science spectra is not an independently measured
background correction. Plot uncertainty and distinguish conditional GP means
from new-noise posterior predictions. A constant GP component is
indistinguishable from $c_p$: show projected variations or differences,
not an absolute recovered continuum. The stellar-frame low-order spectral
basis restricts the GP differently from the Doppler-shifting planet, but
does not guarantee separation. Long timescales and extrapolation across the
science/eclipse gap can remain prior-sensitive.

\textbf{Code and records.}
\path{src/escape/time_gp.py} implements sufficient statistics, Mat\'ern
covariance, marginalized likelihood and conditional coefficients.
Experiments live under \texttt{examples/}: \texttt{gp\_halpha.py} and
\texttt{gp\_controls.py}.


\path{mceg/contracts/escape.time_gp.yaml} records equations, source Bindings,
execution evidence and scoped Assessments. Historical contracts and
full-atmosphere readiness remain unchanged. MCEG is not a runtime dependency;
NumPyro is an optional example dependency.

\small References:
\href{https://arxiv.org/html/2606.23899v1#S2}{Zhang et al. (2026), section II};
\href{https://arxiv.org/abs/1109.3251}{Gibson et al. (2012), GP systematics};
\href{https://num.pyro.ai/en/stable/mcmc.html}{NumPyro NUTS documentation}.
\end{document}
